\chapter[致谢]{致\quad 谢}\chaptermark{致\quad 谢}% syntax: \chapter[目录]{标题}\chaptermark{页眉}
%\thispagestyle{noheaderstyle}% 如果需要移除当前页的页眉
%\pagestyle{noheaderstyle}% 如果需要移除整章的页眉
时间过得真快，仿佛本科毕业还在昨天，今天却又该说告别了。
不过我很荣幸能够在硕士毕业论文的结尾，向我一路走来帮助过我的人说一声最真诚的感谢。

首先，在硕士阶段，我最想感谢的人是我的导师--杨晓雷研究员。
三年来杨老师给予了我太多帮助和指导。
我记得在22年底春节期间，杨老师为了我在来年转博前有发表的文章，几乎每天都在办公室帮我修改论文，我十分感动。
我写的文章初稿，从单词语法到文章结构乃至文章内容，老师几乎对每个句子都做了修改，并且告诉我这样做的原因。
在论文返修的阶段，杨老师也是经常通宵给我改文章，
就是在这样一来一回的反复修改锤炼当中，我的文章撰写能力有了相当的进步。
尤其是在心态上，我不再觉得写论文是一件难事，我的第二篇论文初稿就是我在七天之内撰写完毕的，这完全得益于老师的指导。
%如果没有杨老师的帮助，我是不可能如此迅速地将这些文章发表出来的。
杨老师在科研上严谨负责的态度与始终如一的热情一直影响着我。
我记得第一篇文章返修的时候，我在文章发展的解析模型无法准确预测高湍流来流的情况，这导致审稿人产生了疑问。
我当时只是觉得这完全是因为高湍流来流是一个特例情况，不仅仅是我的模型，文献中的模型也无法准确预测，所以想糊弄过去。
但老师对我说，审稿人的职责就是挑我毛病的不是帮我糊弄文章的，无法预测高湍流来流的情况就是说明模型预测是有问题的，很有可能是我模型没有应用准确，做科学问题一定要细致，不能把问题放过去。
我很庆幸，当时老师对我的教导，我现在慢慢变成了对论文和工作上”斤斤计较“的人。
%
杨老师总是鼓励我，很尊重我。
三年来，杨老师几乎没有对我说过一句语气重的话，总是微笑着鼓励我继续向前走。
即使我偶尔会因为自己的迷糊犯错误，老师也是以”教“为主而”训“则很少。
每次来杨老师办公室讨论问题，他也总是起身回应我，看我流汗时还很体贴地帮我打开空调。
我随着年岁的增长，自己的事情慢慢变多，难免有时候因为工作进展不顺利，导致情绪上有很大的波动。
每当这时候，我想起来老师平时润物细无声的教导，在敬佩他的同时也会平和自己的心态。
%
最后，谢谢杨老师对我一直以来的关心和支持。
去年九月份，我因为申博不顺利，一度想要放弃继续读博的想法。
杨老师在安慰我的同时也帮助我做了申博的攻略。
尤其支持我向外面多走走，如果我没有申请到也会帮助我留在组里继续工作。
同时老师也在心理上为我做建设，让我明白并不是所有的事情都能顺风顺水，读博更是如此，不能因为眼前的一点困难就轻言放弃。
在跟随杨老师学习的这段时间里，我不仅收获了大量的专业知识与学习方法，同时塑造了良好的价值观和学到了为人处世的道理，十分感谢老师的栽培。

我也十分感谢课题组内的王士召老师，杨子轩老师，晋国栋老师对我的尽心帮助和指导。
谢谢李\hbox{\scalebox{0.5} [0.5]{}\lower-1.0ex\hbox{\scalebox{1}[0.45]{明}}\lower.1ex\hbox{\kern-1em \scalebox{1}[0.45]{空}}}斌师兄对我的鼓励和关心，我在程序和工作上的许多问题多亏有师兄的支持才得以解决，师兄经常推荐我，让我有机会学习一些新知识，这对我十分有益。
谢谢周志登师兄对我的帮助，每次在答辩前夕，多亏有师兄指出我PPT的不足并给出相关建议，我才得以顺利地完成汇报。
董国丹师姐和秦建华师兄也常常对我进行耐心的指导，谢谢你们。
还有李世隆师兄和陈丹阳师兄，我会一直怀念和你们在篮球场一起并肩作战的时光。
作为捧哏的刘晓豪师兄我是不会忘记的，为人亲和，关心师弟，和你在一起聊天我总是能开怀大笑。
感谢办公室238c，240E，502内课题组的全体师兄，你们对我的帮助，让我一次又一次度过了难关。

写到这儿，我还要谢谢我在力学所的这段时间遇到的朋友们。
寇军大哥是我不可能忘了的人，带我晨练，带我在宿舍涮羊肉，还有一起爬长城，赏春色，我们相互帮助但也相互吐槽，有太多回忆与感动。
另外澄清一下，虽然军哥打球不传球，但据说这是军哥家乡的风俗，我个人表示一万个理解，并且早已忘了这件事情。
家琦还有文进，我很高兴与你们相识，从雁栖湖的篮球场能打到中关村的篮球场。
刚华是一个很好的人，我们在很多话题和领域都能聊到一起去，谢谢你三年来对我的鼓励和安慰，每次和你的约定你从来没有爽约过，我一直很感动。 
还有好运，和你们在活动室打台球，一起学习进步的时光总是令人怀念。
还有俊淇，孙俊，友华，亚冬，明璨，我的室友-晨冬，等等诸位，谢谢你们一直在我身边，对我总是支持，理解还有关心。
谢谢我的学弟-俊鹏在北京的这段时间对我的陪伴，对我，你总是不怕麻烦，有求必应，希望你之后能学业顺利。
正是因为有你们这些朋友的存在，我在工作之余才能有丰富多彩的生活。

在致谢的最后，我要谢谢我最亲近的人一直以来对我的陪伴和奉献。
首先我要感谢我的父母，你们总是理解也支持着我的选择，并一直对我感到骄傲，这带给我了极大的安慰。
你们将一切力所能及的都给了我，如果没有你们为我遮风挡雨就没有我现在的一切，所以在我眼里你们就是最伟大的人。
谢谢你们能够一直陪在我的身边，我才得以有足够的勇气继续前进。
感谢姥姥和姥爷还有小姨对我的牵挂和照顾，我会一直记得你们对我的恩情。
还有我的诸多亲友，我也十分感谢你们，是你们组成了我对家乡的思念，无论我离家多远，我都不会忘了回来看看你们。
我还要感谢我的女朋友--郭子硕同学，你对我的体贴和包容我看在眼里，你对我的付出我也记在心上。
不知不觉我们已经完成了三年的异地恋挑战，往后肯定还有更多的挑战，和你在一起我有信心。


正是因为有大家的陪伴和帮助，我一直觉得我是一个很幸福的人。
谨以此篇致谢献给帮助我的所有人。
